\chapter{Conclusiones y trabajo futuro}
\label{cap:conclusiones}

\section{Síntesis y aportación del trabajo}
\hypertarget{RE3.2}{}\label{RE3.2}

El trabajo entrega una cadena completa y cerrada, que empieza en un conjunto
público de formas de onda y termina en una señal de línea en un conector de
3,5~mm. Entre esos dos extremos hay un autoencoder variacional entrenado en el
ordenador, una rejilla de 256 tablas horneada a partir
de su decodificador, tres microcontroladores con una responsabilidad única cada
uno, dos enlaces de datos verificados muestra a muestra y una carcasa que lo
alberga todo. El reparto que hace posible ese encadenamiento se describe en el
apartado~\ref{sec:arquitectura}. Lo que ese conjunto demuestra es que un espacio
tímbrico aprendido se puede recorrer con el dedo sobre hardware de catálogo, con
83,82~euros de materiales según el Apéndice~\ref{anx:presupuesto}, sin un
ordenador delante y sin ningún circuito integrado hecho a medida.

Conviene dejar dicho con precisión qué papel juega aquí el aprendizaje
automático, porque es donde más fácil resulta exagerar. En el instrumento no se
ejecuta ninguna red neuronal. El codificador solo intervino durante el
entrenamiento y el decodificador se ejecutó fuera de línea para hornear la
rejilla, de modo que en marcha solo hay consulta de tabla e interpolación. La
inteligencia está en el contenido de las tablas y no en el código que las lee,
ya que el hecho de que celdas vecinas suenen emparentadas y de que cruzar el
plano produzca una transición tímbrica coherente es producto del latente
aprendido y no se obtiene a mano. Precalcular fue una estrategia de despliegue y
no una renuncia, y el motivo está en el tiempo de cálculo y no en el espacio,
tal como se razonó en el apartado~\ref{sec:alternativas}. El tamaño de memoria, de hecho,
no distinguía entre las dos opciones. La rejilla ocupa 512~kB
(apartado~\ref{sec:rejilla}) y el decodificador, con sus cerca de 592.000
parámetros, ocuparía del orden de 590~kB cuantizado a ocho bits. Lo que la
rejilla compra es un coste de acceso constante y conocido. Lo que paga a cambio
es una resolución del plano congelada en el momento de hornearla.

\section{Cumplimiento de los objetivos}
\hypertarget{RE3.1}{}\label{RE3.1}

Los tres primeros objetivos específicos corresponden a la fase de preparación y
quedan cumplidos. El conjunto de datos homogéneo se construyó resolviendo el
remuestreo, la normalización de amplitud y la alineación de fase
(apartado~\ref{sec:preprocesado}). El autoencoder variacional comprime cada onda
a una coordenada de dos dimensiones, y su plano resultó continuo y bien
aprovechado (apartado~\ref{sec:latente}), con la diversidad confirmada después
sobre la propia rejilla (apartado~\ref{sec:demo}). Y esa rejilla se horneó y se
exportó en un formato que el firmware compila directamente
(apartado~\ref{sec:rejilla}).

Los tres siguientes son de firmware y también quedan cumplidos. El motor
sustituye la tabla en reproducción sin discontinuidad audible, con una
transición cuyo rizado de pico se midió en 0,34~dB
(apartado~\ref{sec:crossfade}). La comunicación entre las tres placas quedó
definida y verificada, con la interpolación equivalente a su referencia dentro
de 1~\gls{lsb} y el transporte por \gls{spi} sin una sola muestra alterada
(apartado~\ref{sec:validacion-datos}). El control de nota y dinámica por
\gls{midi} funciona, con las 128 notas dentro de tolerancia y el la central en
440~Hz sobre el montaje definitivo, y el panel analógico recorre el fondo de
escala completo en todos sus canales (apartado~\ref{sec:tensiones}). Este último
objetivo se cumple con una salvedad que no conviene disimular, y es que la
realimentación visual de los mandos no llegó a funcionar sobre el aparato
montado (apartado~\ref{sec:incidencias}).

Quedan el objetivo de validar cada etapa contra una referencia calculada en el
ordenador, que atraviesa todo el desarrollo y se recoge reunido en el
capítulo~\ref{cap:pruebas}, y el de integrar el conjunto en un prototipo físico
alimentado desde una fuente única, que consume 0,42~A y entrega 198,7~mV
eficaces de nivel de línea (apartados~\ref{sec:consumo} y~\ref{sec:audio}).
También este último lleva salvedad, y es que el procesador de audio trabaja
medio voltio por debajo de la tensión de entrada mínima que garantiza su
fabricante, con la causa identificada y la corrección propuesta
(apartado~\ref{sec:balance}). Con esas dos salvedades sobre la mesa, el objetivo
general se considera alcanzado. Hay un instrumento autónomo que permite explorar
un espacio tímbrico aprendido con una pantalla táctil y una entrada \gls{midi}
estándar.

\section{Aprendizajes del proyecto}

Por encima de cualquier resultado concreto, el trabajo ha permitido aplicar
sobre un mismo objeto contenidos que en el grado se estudian por separado. El
aprendizaje automático entrena el modelo, el procesado de señal prepara las
ondas y las reproduce, los sistemas embebidos reparten la carga entre tres
procesadores, la electrónica analógica resuelve el panel y la alimentación, los
protocolos de comunicación unen las placas y el diseño mecánico las encierra en
algo que cabe en las manos. Que todas esas materias tuvieran que encajar en un
único aparato, y no cada una en su práctica, ha sido lo más satisfactorio del
proyecto, y es la razón por la que se eligió este tema y no otro más acotado.

\section{Trabajo futuro}
\hypertarget{RE3.2b}{}\label{RE3.2b}

Las líneas de continuación salen todas de una limitación conocida, y conviene
enunciarlas con ella delante. La primera es la que el capítulo~\ref{cap:diseno}
dejó abierta a propósito. El decodificador puede llegar a ejecutarse a bordo del
ESP32-S3, y el reparto elegido lo permite sin tocar el firmware del procesador
de audio, que seguiría recibiendo exactamente lo mismo
(apartado~\ref{sec:alternativas}). Antes de intentarlo conviene saber si merece
la pena, y esa comprobación no necesita hardware. Como el decodificador no es
lineal, interpolar entre dos de sus salidas no equivale a decodificar el punto
latente intermedio, de modo que comparar ambas cosas en puntos intermedios daría
la cifra de error de aproximación que hoy no existe, y con ella se sabría si una
rejilla de $16 \times 16$ basta. La segunda línea es la limitación de banda de
las tablas, que quedó fuera del alcance porque corregirla multiplica el volumen
de datos embarcado por el número de versiones filtradas que haya que almacenar
(apartado~\ref{sec:parametros}), y que se resolvería eligiendo la adecuada según
el registro de la nota. La tercera es la fidelidad del cuello de botella de dos
dimensiones, que redondea los flancos verticales
(apartado~\ref{sec:reconstruccion}) y cuya ampliación obligaría a replantear el
control táctil, que es precisamente lo que fijó esas dos dimensiones.

Antes que ninguna de ellas está la revisión del propio aparato. Las tres
incidencias que la verificación dejó abiertas siguen abiertas
(apartado~\ref{sec:incidencias}). Falta además levantar con
osciloscopio y analizador lo que un multímetro no alcanza, esto es, la
distorsión, la respuesta en frecuencia y el pico de corriente del arranque
(apartado~\ref{sec:balance}). La polifonía, descartada desde el principio, solo
tiene sentido plantearla cuando todo lo anterior esté resuelto.

El prototipo es lo más perecedero de todo esto. Lo que permanece no está dentro
de la caja, sino dentro de mi cabeza. La experiencia de haberlo hecho, de haberlo 
pensado y de haberlo escrito es lo que queda, y es lo que me permitirá abordar con
más seguridad los próximos proyectos. Proyectos serán más ambiciosos, proyectos que espero que me dejen
con la misma sensación de haber aprendido algo que no estaba en ningún libro ni en ninguna clase. 

